\subsection{消息与参与方互异关系}
\label{subsec:relation-discussion}

参与方、发送实例和消息三个维度分别描述同批条目的主体归属、发送发生和消息内容。沿用第\ref{subsec:configurations-properties}节的记号，机器结果直接支持$M_\tau\nRightarrow P_\tau$：\path{same_party_different_messages_batch_exists}给出同一参与方的不同消息在同批被接受的完整可达执行。与之不同，$P_\tau\Rightarrow M_\tau\Leftrightarrow I_\tau$来自三个模型共同采用的新鲜消息生成和完整元组精确来源匹配规则，是人工迹语义推导，而非新增的Tamarin lemma或证明器结果。该关系不能推广为任意协议中的一般性质，也不意味着三个性质具有相同语义或三个比较维度彼此独立。

发送来源对应性和批内来源单射性约束发送—接受事件关系：前者要求每个\receiveraccept 都有更早且与完整$(A,\oid,m)$匹配的\sendaction，后者限制同一匹配发送来源在一个$(\mathit{bid},\mathit{rst})$中对应的接受事件数。两者均不比较不同来源的参与方投影，批内来源单射性也不同于Lowe认证层次中的单射一致性。因此，两个条目可以分别满足来源要求并各自接受一次，同时仍具有相同的参与方投影。

第\ref{subsec:message-results}节的机器见证将上述维度差异落实为可达行为：两个条目具有不同发送实例、不同消息和各自的匹配发送来源，但参与方投影均为$A$，并在同批产生接受事件。该执行的意义不在于重述“不同消息不等于不同参与方”，而在于证明这种组合能够通过共同的来源匹配、准入和顺序处理生命周期。因此，条目级的消息和来源性质不能替代式\eqref{eq:distinct-party-objective}规定的批次组合关系。

\subsection{批次准入接口的关系维护}
\label{subsec:interface-maintenance}

式\eqref{eq:distinct-party-objective}规定同批条目应保持的关系；参与方互异约束模型在准入处直接比较$A_1$与$A_2$，为该关系提供可执行的目标对照。该建模选择用于验证目标关系及有效批次的可达性，不说明显式比较是唯一实现方式。

该关系可以由调用者、批次构造组件、准入层或更早的数据结构不变量维护，现有模型不验证这些机制的具体实现，也不比较其优劣。若系统需要保持K-Waay原有的批内参与方互异条件，负责组合或准入的受信任边界必须获得适当的参与方表示，并保证条目进入同一批次时该关系成立。这一接口义务既不指定唯一的责任组件，也不要求实现采用固定的条件分支。

参与方表示仍需由具体系统定义。模型中的$A$是抽象协议主体，不自动等同于账户、长期身份、公钥编码、设备、预密钥关联或数据库记录。系统只有建立实现对象到参与方投影的合理映射，才能据此比较批内条目；当前模型既未验证该映射，也未引入新的身份解析机制。

\subsection{适用范围与局限}
\label{subsec:scope-limitations}

机器证据固定为两个槽位。式\eqref{eq:distinct-party-objective}通过成对比较定义有限批次的目标关系，两槽模型捕获最小非平凡相等情形，但没有证明任意批次长度下的完整协议性质。结论仅涉及同一$(\mathit{bid},\mathit{rst})$中的条目组合和接受事件，不提供跨批次、回滚或重启保证。

模型以持久事实\texttt{!Sent(A,oid,m)}表示匹配发送来源。该事实没有实现签名验证、KEM与split-KEM解封装、密钥导出或其他真实认证过程；发送实例和消息的新鲜性也是当前抽象的建模选择。模型不含完整K-Waay安全游戏中的\textsc{Key}或\textsc{Test}接口，直接证据止于批次准入和\receiveraccept，因而不能推出密钥泄露、机密性失败或完整认证密钥交换性质失效。

接受事件后的消费组件、状态安装和会话使用过程同样未被建模，完整K-Waay执行与两槽抽象之间也没有精化定理。因此，重复接受不等于重复状态安装或应用状态破坏，抽象中的可达执行也不能直接解释为具体协议攻击。攻击者控制批次组合属于分析假设；现有证据既不表明真实部署中的攻击者已经获得该能力，也不表明实际实现遗漏了原有条件。

两项拒绝性质只证明相应分支存在可达执行，不构成所有无效输入最终被拒绝的活性结论。其中，参与方拒绝性质尚未把先前的\sendaction 元组绑定到被拒候选元组。2026年9月16日的独立复核确认三个未修改模型在记录环境中的结果状态和证明步数一致，支持当前抽象的可复现性，但不扩大协议级或部署级结论。

